\section{MÔ HÌNH HÓA BÀI TOÁN}
\label{sec:chuong2}

\subsection{Dữ liệu đầu vào và các giả định}

Một đầu vào xác định kho hàng, tập đơn cần phục vụ và các thông số vận hành. Gọi $O$ là tập đơn, $K$ là tập sản phẩm và $P=\{1,\ldots,m\}$ là tập nhân viên, với $n=|O|>0$. Đơn $i\in O$ yêu cầu số lượng $q_{is}$ của sản phẩm $s$; quy ước $q_{is}=0$ nếu đơn không yêu cầu sản phẩm đó. Mỗi sản phẩm có vị trí lấy hàng $\ell_s$, tải của một đơn vị $a_s>0$ và thời gian lấy một đơn vị $h_s\geq0$. Tải của đơn được tính bởi
\begin{equation}
 w_i=\sum_{s\in K}a_sq_{is}.\label{eq:model-order-load}
\end{equation}
Đơn có hạn hoàn thành $d_i\geq0$ và một thứ tự đầu vào. Thứ tự này phục vụ việc gom đơn cơ sở, không phải thời điểm đơn phát sinh.

Thông số vận hành gồm sức chứa xe $Q>0$, vận tốc di chuyển $v>0$, thời gian chuẩn bị mỗi chuyến $b\geq0$ và thời gian thao tác cố định tại mỗi vị trí có yêu cầu lấy hàng $h_{\rm loc}\geq0$. Sức chứa và tải sản phẩm dùng cùng một đơn vị tải quy ước, không mặc định là khối lượng. Dữ liệu mô phỏng dùng mét và phút; dữ liệu Kris giữ quy ước số học của nguồn. Trong từng đầu vào, khoảng cách, vận tốc và thời gian phải nhất quán. Việc quy đổi đơn vị hiển thị không thay thế việc đánh giá nghiệm theo dữ liệu nguồn.

\begin{table}[H]
\centering\small
\begin{tabularx}{\linewidth}{lX}
\toprule
Ký hiệu & Ý nghĩa \\
\midrule
$O,K,P$ & Tập đơn, tập sản phẩm và tập nhân viên. \\
$n,m$ & Số đơn và số nhân viên. \\
$q_{is},a_s,h_s$ & Số lượng yêu cầu, tải và thời gian lấy một đơn vị sản phẩm. \\
$\ell_s,w_i,d_i$ & Vị trí sản phẩm, tải đơn và hạn hoàn thành. \\
$Q,v,b,h_{\rm loc}$ & Sức chứa, vận tốc, thời gian chuẩn bị và thao tác tại vị trí. \\
$B_{p,j},J_p$ & Tập đơn trong chuyến thứ $j$ và số chuyến của nhân viên $p$. \\
$U(B),D(B),\tau(B)$ & Tập vị trí cần lấy hàng, quãng đường và thời lượng chuyến. \\
$S_{p,j},C_{p,j}$ & Thời điểm bắt đầu và kết thúc chuyến. \\
$C_i,T_i$ & Thời điểm hoàn thành và độ trễ của đơn. \\
$D,C_{\max},T,L$ & Tổng quãng đường, makespan, tổng độ trễ và số đơn trễ. \\
$t_{\rm run}$ & Thời gian máy tính chạy để xây dựng và kiểm chứng kế hoạch. \\
$\alpha_D,\alpha_C,\alpha_T$ & Trọng số của ba thành phần trong hàm mục tiêu. \\
\bottomrule
\end{tabularx}
\caption{Các ký hiệu sử dụng trong mô hình bài toán.}\label{tab:model-symbols}
\end{table}

Mô hình sử dụng các giả định sau:
\begin{enumerate}
\item Toàn bộ đơn được biết và sẵn sàng tại thời điểm $0$; không bổ sung đơn trong khi thực hiện kế hoạch.
\item Nhân viên đồng nhất về vận tốc và thời gian thao tác. Mỗi người dùng xe có cùng sức chứa và thực hiện tối đa một chuyến tại một thời điểm.
\item Mỗi đơn được phục vụ trọn vẹn trong đúng một chuyến. Đầu vào phải thỏa $w_i\leq Q$, nên một đơn riêng lẻ không cần bị tách để đáp ứng sức chứa.
\item Mỗi chuyến bắt đầu và kết thúc tại depot. Đơn hoàn thành khi chuyến chứa nó quay về depot.
\item Hạn hoàn thành là mềm: trễ hạn làm tăng chi phí nhưng không làm nghiệm mất tính khả thi.
\item Không xét tắc nghẽn, xung đột di chuyển, thiếu hàng, ca nghỉ hoặc sự thay đổi vận tốc theo tải.
\end{enumerate}
Các giả định xác định phạm vi áp dụng của kết quả. Một nghiệm hợp lệ trong mô hình chưa đủ để bảo đảm đáp ứng mọi điều kiện vận hành của một kho thực tế.

\subsection{Biểu diễn kho hàng bằng đồ thị}

Kho được biểu diễn bằng đồ thị vô hướng liên thông $G=(V,E)$ có trọng số cạnh không âm. Tập $V$ gồm depot, vị trí sản phẩm và các điểm trung chuyển; cạnh $(u,v)\in E$ biểu diễn một đoạn lối đi được phép sử dụng. Gọi $c_{uv}$ là chiều dài của đoạn đó và $d_0\in V$ là depot. Tọa độ của đỉnh dùng để thể hiện bố trí kho; quãng đường được tính từ các cạnh, không dùng khoảng cách đường thẳng xuyên qua kệ.

Khoảng cách giữa hai vị trí là độ dài đường đi ngắn nhất trên đồ thị:
\begin{equation}
 \delta(u,v)=\min_{\pi:u\leadsto v}\sum_{(x,y)\in\pi}c_{xy}.\label{eq:model-graph-distance}
\end{equation}
Với trọng số không âm, chương trình dùng Dijkstra để tính khoảng cách và lưu thông tin truy vết đường đi. Do đồ thị liên thông, mọi vị trí cần phục vụ đều có thể tiếp cận từ depot.

Với chuyến chứa tập đơn $B$, tập vị trí có yêu cầu lấy hàng là
\begin{equation}
 U(B)=\bigcup_{i\in B}\{\ell_s:q_{is}>0\}.\label{eq:model-required-locations}
\end{equation}
Nhiều sản phẩm hoặc nhiều đơn có thể dùng chung một vị trí; vị trí đó chỉ xuất hiện một lần trong $U(B)$. Nếu thứ tự điểm dừng là $(u_0,u_1,\ldots,u_r,u_{r+1})$, trong đó $u_0=u_{r+1}=d_0$ và các điểm lấy hàng ngoài depot nằm giữa hai đầu, quãng đường chuyến được tính bởi
\begin{equation}
 D(B)=\sum_{t=0}^{r}\delta(u_t,u_{t+1}).\label{eq:model-batch-distance}
\end{equation}
Từng đoạn nối điểm dừng sau đó được mở rộng thành đường đi trên các cạnh của kho. Vì vậy, dãy điểm dừng phục vụ và đường đi vật lý là hai cách biểu diễn khác nhau: đường đi vật lý có thể qua một đỉnh nhiều lần dù vị trí lấy hàng chỉ được liệt kê một lần trong dãy điểm dừng.

\subsection{Biểu diễn phương án gom đơn, tuyến đi và lịch thực hiện}

Ở mức lập kế hoạch, một phương án cần xác định các đơn thuộc từng chuyến, tuyến phục vụ của chuyến, nhân viên phụ trách và thứ tự thực hiện. Trong cài đặt, cấu trúc nghiệm chính là danh sách chuyến có thứ tự của từng nhân viên:
\begin{equation}
 X=(\sigma_1,\ldots,\sigma_m),\qquad
 \sigma_p=(B_{p,1},\ldots,B_{p,J_p}).\label{eq:model-plan}
\end{equation}
Một nhân viên không được giao việc có $J_p=0$. Tập đơn của chuyến quyết định $U(B)$; thủ tục định tuyến dựng tuyến trên đồ thị, rồi bộ đánh giá tính thời lượng và lịch theo thứ tự chuyến.

Như vậy, cấu trúc nghiệm trong chương trình lưu trực tiếp phân hoạch đơn và thứ tự chuyến, còn tuyến và thời điểm được suy ra bằng bộ giải mã. Đây là cách áp dụng heuristic cho bài toán tích hợp; chương trình không tìm kiếm độc lập trên toàn bộ các tuyến và các thời điểm bắt đầu có thể có. Quy tắc định tuyến cụ thể của từng thuật toán được trình bày ở Chương 3.

Thay đổi một đơn giữa hai chuyến có thể làm thay đổi cả tuyến và thời lượng. Đổi thứ tự hai chuyến hoặc chuyển chuyến sang nhân viên khác có thể làm thay đổi nhiều thời điểm hoàn thành phía sau. Vì vậy, sau khi thay đổi nghiệm, cần tính lại các chỉ tiêu liên quan trước khi so sánh chất lượng.

\subsection{Các ràng buộc của phương án khả thi}

Gọi $\mathcal F$ là tập các cấu trúc nghiệm hoàn chỉnh thỏa các điều kiện dưới đây. Trước hết, mọi đơn phải được phục vụ đúng một lần:
\begin{align}
 \bigcup_{p=1}^{m}\bigcup_{j=1}^{J_p}B_{p,j}&=O,\label{eq:model-cover}\\
 B_{p,j}\cap B_{q,k}&=\varnothing\quad\text{khi }(p,j)\neq(q,k).\label{eq:model-disjoint}
\end{align}
Mỗi chuyến phải chứa ít nhất một đơn và tuân thủ sức chứa:
\begin{equation}
 B_{p,j}\neq\varnothing,\qquad \sum_{i\in B_{p,j}}w_i\leq Q.\label{eq:model-capacity}
\end{equation}
Nghiệm chỉ sử dụng các đơn có trong đầu vào. Trạng thái trung gian tạm thiếu đơn trong quá trình tìm kiếm không được xem là một nghiệm đầu ra hoàn chỉnh.

Tuyến của mỗi chuyến phải bắt đầu và kết thúc tại depot, phục vụ đủ $U(B)$ và chỉ di chuyển trên các cạnh hợp lệ. Nhân viên thực hiện các chuyến tuần tự; chuyến sau không bắt đầu trước khi chuyến trước kết thúc:
\begin{equation}
 S_{p,1}\geq0,\qquad S_{p,j+1}\geq C_{p,j}.\label{eq:model-no-overlap}
\end{equation}
Trong cài đặt, chuyến đầu bắt đầu tại $0$, chuyến sau bắt đầu ngay khi chuyến trước kết thúc. Với toàn bộ đơn sẵn sàng tại $0$ và không có ràng buộc giờ mở cửa hoặc thời điểm sớm nhất, loại bỏ thời gian chờ trong một thứ tự chuyến cố định không làm tăng quãng đường, makespan hoặc độ trễ.

Không đặt ràng buộc $C_i\leq d_i$ vì mô hình cho phép trễ hạn. Cần phân biệt kiểm chứng khả thi với đánh giá chất lượng: nghiệm đủ đơn và đúng sức chứa vẫn có thể có độ trễ lớn hoặc giá trị mục tiêu cao. Chương trình kiểm tra đầu vào, cấu trúc nghiệm và đối chiếu tuyến, thời gian, chỉ tiêu trong kết quả xuất ra; các phép so sánh số thực có sử dụng dung sai.

\subsection{Hàm mục tiêu và các chỉ tiêu đánh giá}

Thời lượng chuyến gồm chuẩn bị, di chuyển, thao tác tại các vị trí cần lấy hàng và lấy từng đơn vị sản phẩm:
\begin{equation}
 \tau(B)=b+\frac{D(B)}{v}+h_{\rm loc}|U(B)|
 +\sum_{i\in B}\sum_{s\in K}h_sq_{is}.\label{eq:model-duration}
\end{equation}
Thời gian thao tác tại vị trí được tính một lần cho mỗi vị trí có yêu cầu trong chuyến, kể cả khi nhiều đơn dùng chung vị trí đó. Thời gian lấy sản phẩm vẫn tính theo tổng số lượng. Việc đường đi vật lý qua một vị trí nhiều lần không làm tự động tăng số lần tính thao tác tại vị trí trong mô hình.

Lịch được tính liên tiếp theo từng nhân viên:
\begin{align}
 C_{p,0}&=0,\label{eq:model-schedule-origin}\\
 S_{p,j}&=C_{p,j-1},\qquad
 C_{p,j}=S_{p,j}+\tau(B_{p,j}).\label{eq:model-schedule}
\end{align}
Mọi đơn thuộc một chuyến có cùng thời điểm hoàn thành, bằng thời điểm chuyến quay về depot:
\begin{equation}
 C_i=C_{p,j}\quad\text{với }i\in B_{p,j},\qquad
 T_i=\max(0,C_i-d_i).\label{eq:model-order-completion}
\end{equation}

Ba thành phần chính là tổng quãng đường, thời gian hoàn tất toàn bộ công việc và tổng độ trễ:
\begin{align}
 D(X)&=\sum_{p=1}^{m}\sum_{j=1}^{J_p}D(B_{p,j}),\label{eq:model-total-distance}\\
 C_{\max}(X)&=\max_{p\in P}C_{p,J_p},\label{eq:model-makespan}\\
 T(X)&=\sum_{i\in O}T_i.\label{eq:model-total-tardiness}
\end{align}
Quy ước $C_{p,J_p}=0$ cho nhân viên không có chuyến. Makespan là thời điểm người cuối cùng kết thúc, không phải tổng thời gian làm việc của các nhân viên. Tổng độ trễ cộng phần vượt hạn của từng đơn, nên một chuyến có nhiều đơn trễ có thể đóng góp nhiều lần vào $T$.

Các đại lượng $S_{p,j}$, $C_{p,j}$ và $C_{\max}$ thuộc lịch vận hành, được tính từ tuyến và thời gian thao tác; với dữ liệu mô phỏng, chúng có đơn vị phút. Ngược lại, thời gian tính toán $t_{\rm run}$ được đo bằng đồng hồ máy tính theo giây, bao gồm tiền xử lý, khởi tạo, tìm kiếm và kiểm chứng của lượt giải. Đây là chỉ tiêu đánh giá chi phí tính toán, không phải thời lượng chuyến và không nằm trong hàm mục tiêu $F$. Khi phân tích kết quả, thời gian tính toán và makespan được báo cáo riêng.

Chương trình ghi thêm số đơn trễ và tỷ lệ đơn đúng hạn:
\begin{equation}
 L(X)=\sum_{i\in O}\mathbf1\{T_i>\varepsilon\},\qquad
 R_{\rm on}(X)=1-\frac{L(X)}{n},\qquad \varepsilon=10^{-9}.\label{eq:model-service-metrics}
\end{equation}
$L$ là số đếm, còn $T$ là tổng thời gian; hai chỉ tiêu phản ánh các khía cạnh khác nhau. Dung sai $\varepsilon$ được dùng trong cài đặt để nhận diện trễ khi tính bằng số thực. Các công thức thời lượng, lịch và chỉ tiêu trong mục này tương ứng với bộ đánh giá của chương trình.

\subsection{Chuẩn hóa các thành phần và lựa chọn trọng số}

Để kết hợp các đại lượng có đơn vị và thang giá trị khác nhau, đề tài dùng nghiệm First-Come, First-Served Batching (B0) làm mốc trên từng đầu vào. Gọi $D_0$ và $C_0$ là tổng quãng đường và makespan của nghiệm này, tính với quy tắc định tuyến cơ sở. Các thang tham chiếu trong cài đặt là
\begin{equation}
 D_{\rm ref}=\max(D_0,1),\qquad
 C_{\rm ref}=\max(C_0,1),\qquad
 T_{\rm ref}=nC_{\rm ref}.\label{eq:model-normalization}
\end{equation}
Giá trị $1$ là sàn số học trong quy ước đơn vị đầu vào để tránh mẫu số bằng không. Thang độ trễ không lấy từ $T(B0)$: ngay cả khi B0 không có đơn trễ, $nC_{\rm ref}$ vẫn dương. Các thang này được tính một lần từ B0 và giữ cố định khi đánh giá các phương pháp, seed tìm kiếm và cấu hình trọng số trên cùng đầu vào.

Hàm mục tiêu cần tối thiểu hóa là
\begin{equation}
 \min_{X\in\mathcal F}F(X),\qquad
 F(X)=\alpha_D\frac{D(X)}{D_{\rm ref}}
 +\alpha_C\frac{C_{\max}(X)}{C_{\rm ref}}
 +\alpha_T\frac{T(X)}{T_{\rm ref}},\label{eq:model-objective}
\end{equation}
với $\alpha_D,\alpha_C,\alpha_T>0$ và $\alpha_D+\alpha_C+\alpha_T=1$. Tuyến và lịch trong $F(X)$ được giải mã theo quy tắc của phương pháp đang xét. Biểu thức tối thiểu hóa xác định mục tiêu của bộ giải, không khẳng định các heuristic tìm được nghiệm tối ưu của tập khả thi.

Cấu hình tham chiếu dùng $(\alpha_D,\alpha_C,\alpha_T)=(1/3,1/3,1/3)$. Trọng số bằng nhau sau chuẩn hóa là quy ước đối chiếu, chưa được hiệu chuẩn từ chi phí kinh tế hoặc mức phục vụ của một kho thực tế. Thành phần $T/(nC_{\rm ref})$ chính là độ trễ trung bình mỗi đơn chia cho makespan tham chiếu. Số đơn trễ $L$ không nằm trực tiếp trong $F$.

Giá trị $F$ thấp hơn biểu thị phương án tốt hơn theo trọng số đã chọn. Tuy nhiên, giảm $F$ không bảo đảm cả ba thành phần đều giảm, và tăng trọng số độ trễ không đồng nghĩa trực tiếp với tối thiểu hóa số đơn trễ. Chỉ đối chiếu điểm $F$ khi dùng cùng đầu vào, cùng thang tham chiếu và cùng trọng số; giữa các cấu hình trọng số, cần xem thêm $D$, $C_{\max}$ và $T$. Các thang tham chiếu được xác định riêng cho từng đầu vào, nên không diễn giải chênh lệch điểm giữa hai bài khác nhau như chênh lệch chi phí vận hành tuyệt đối.

\subsection{Ví dụ tính toán minh họa}

Xét một kho nhỏ có bốn đơn và hai nhân viên. Depot cùng bốn vị trí sản phẩm nằm trên một lối đi như \figref{fig:model-example}. Mỗi cạnh có thể đi theo hai chiều; số ghi trên cạnh là chiều dài thực, còn sơ đồ không vẽ theo tỷ lệ.

\begin{figure}[H]
\centering
\begin{tikzpicture}[x=1cm,y=1cm,every node/.style={font=\small}]
\node[draw,rounded corners=3pt,inner sep=5pt] (depot) at (0,0) {Depot};
\node[draw,inner sep=5pt] (a) at (2.4,0) {$A$};
\node[draw,inner sep=5pt] (b) at (4.3,0) {$B$};
\node[draw,inner sep=5pt] (c) at (6.8,0) {$C$};
\node[draw,inner sep=5pt] (d) at (8.7,0) {$D$};
\draw (depot)--node[above]{2 m}(a);
\draw (a)--node[above]{1 m}(b);
\draw (b)--node[above]{5 m}(c);
\draw (c)--node[above]{1 m}(d);
\node[below=0.25cm of a] {$O_1$};
\node[below=0.25cm of b] {$O_3$};
\node[below=0.25cm of c] {$O_2$};
\node[below=0.25cm of d] {$O_4$};
\end{tikzpicture}
\caption{Đồ thị lối đi của ví dụ bốn đơn.}\label{fig:model-example}
\end{figure}

Mỗi đơn yêu cầu một đơn vị sản phẩm, có tải $1$ và thời gian lấy $1$ phút. Sức chứa xe $Q=2$, vận tốc $v=2$ mét/phút, thời gian chuẩn bị $b=1$ phút và thời gian thao tác tại vị trí $h_{\rm loc}=0$. Thứ tự đầu vào là $O_1,O_2,O_3,O_4$; hạn và vị trí của từng đơn được cho trong \tabref{tab:example-orders}.

\begin{table}[H]
\centering\small
\begin{tabular}{llrr}
\toprule
Đơn & Vị trí sản phẩm & Tải & Hạn, phút \\
\midrule
$O_1$ & $A$ & 1 & 7 \\
$O_2$ & $C$ & 1 & 12 \\
$O_3$ & $B$ & 1 & 7 \\
$O_4$ & $D$ & 1 & 12 \\
\bottomrule
\end{tabular}
\caption{Đơn hàng trong ví dụ minh họa.}\label{tab:example-orders}
\end{table}

B0 gom theo thứ tự thành hai chuyến $\{O_1,O_2\}$ và $\{O_3,O_4\}$, giao cho hai nhân viên và cùng bắt đầu tại $0$. Chuyến thứ nhất đi đến $C$ rồi quay về, với điểm dừng $d_0\to A\to C\to d_0$ và quãng đường $2+6+8=16$ mét. Theo \eqref{eq:model-duration}, thời lượng là $1+16/2+2=11$ phút. Chuyến thứ hai có quãng đường $3+6+9=18$ mét và thời lượng $1+18/2+2=12$ phút.

Vì đơn hoàn thành khi chuyến quay về, $O_1$ trễ $11-7=4$ phút và $O_3$ trễ $12-7=5$ phút; hai đơn còn lại đúng hạn. Nghiệm B0 vẫn khả thi vì mỗi chuyến có tải $2$ và đủ bốn đơn. Các chỉ tiêu là $D=34$, $C_{\max}=12$, $T=9$, $L=2$ và $R_{\rm on}=50\%$. Thang chuẩn hóa là $D_{\rm ref}=34$, $C_{\rm ref}=12$ và $T_{\rm ref}=4\times12=48$.

Với trọng số bằng nhau, điểm của B0 được tính bởi
\begin{equation}
 F(B0)=\frac13\left(\frac{34}{34}+\frac{12}{12}+\frac9{48}\right)
 =\frac{35}{48}\approx0{,}729167.\label{eq:model-example-b0}
\end{equation}

Trên cùng đầu vào, Greedy Batching (B2) tạo hai chuyến $\{O_1,O_3\}$ và $\{O_2,O_4\}$ trong cài đặt hiện tại. Chuyến gần depot có quãng đường $2+1+3=6$ mét và thời lượng $1+6/2+2=6$ phút. Chuyến xa có quãng đường $8+1+9=18$ mét và thời lượng $12$ phút. Mỗi nhân viên nhận một chuyến; cả bốn đơn đều đúng hạn.

Để minh họa riêng ảnh hưởng của lập lịch, xét phương án $X$ giữ nguyên hai chuyến và tuyến của B2 nhưng giao cả hai cho nhân viên thứ nhất, thực hiện chuyến gần trước. Nhân viên thứ hai không nhận chuyến. Chuyến đầu bắt đầu tại $0$ và kết thúc tại $6$; chuyến sau bắt đầu tại $6$ và kết thúc tại $18$ phút. Hai đơn gần depot vẫn đúng hạn, còn $O_2,O_4$ cùng trễ $18-12=6$ phút. Vì vậy $D=24$, $C_{\max}=18$, $T=12$ và $L=2$. Đây là phương án phân công khác dùng để minh họa, không phải đầu ra của B2.

\begin{table}[H]
\centering\small
\begin{tabular}{lrrrrr}
\toprule
Phương án & $D$, m & $C_{\max}$, phút & $T$, phút & $L$ & $F$ \\
\midrule
B0 & 34 & 12 & 9 & 2 & 0,729167 \\
B2 & 24 & 12 & 0 & 0 & 0,568627 \\
$X$, một nhân viên & 24 & 18 & 12 & 2 & 0,818627 \\
\bottomrule
\end{tabular}
\caption{So sánh các chỉ tiêu của ba phương án trong ví dụ.}\label{tab:example-comparison}
\end{table}

Dùng cùng thang chuẩn hóa của B0, điểm của B2 là
\begin{equation}
 F(B2)=\frac13\left(\frac{24}{34}+\frac{12}{12}+\frac0{48}\right)
 =\frac{29}{51}\approx0{,}568627.\label{eq:model-example-b2}
\end{equation}
Theo \tabref{tab:example-comparison}, B2 giảm $10$ mét, giữ nguyên makespan và loại độ trễ trong ví dụ. Mức giảm mục tiêu so với B0 là
\begin{equation}
 100\frac{F(B0)-F(B2)}{F(B0)}\approx22{,}02\%.\label{eq:model-example-improvement}
\end{equation}
Dùng cùng thang chuẩn hóa của B0, điểm của phương án $X$ là
\begin{equation}
 F(X)=\frac13\left(\frac{24}{34}+\frac{18}{12}+\frac{12}{48}\right)
 =\frac{167}{204}\approx0{,}818627.\label{eq:model-example-sequential}
\end{equation}
So với B2, phương án $X$ giữ nguyên quãng đường nhưng có makespan và độ trễ lớn hơn do dồn hai chuyến vào một nhân viên. Ví dụ này làm rõ rằng đánh giá cách gom đơn và tuyến đi cần đi cùng đánh giá phân công và lịch thực hiện.

Ví dụ cho thấy cách tính các nghiệm khả thi và các thành phần mục tiêu. Cả ba phương án đã được đánh giá và kiểm chứng chỉ tiêu bằng chương trình, nhưng không được cộng vào benchmark hoặc dùng làm bằng chứng về ưu thế trên toàn bộ dữ liệu. Ví dụ cũng không chứng minh B2 đạt tối ưu toàn cục; cơ chế tạo hai cách gom được giải thích tiếp ở Chương 3.
